\documentclass[twoside]{article}
\usepackage[utf8]{inputenc}

% The package supports other formats as well, but this template
% is primarily checked and updated for examinations at TUM BGDM
%\usepackage[uebung, answers]{tumbgdm}

\usepackage[uebung]{tumbgdm}
\usepackage{pythonhighlight}

\Professor{Prof. Dr. Martin Werner}
\Universitaet{Technical University of Munich}
\Institut{School of Engineering and Design}
\Lehrstuhl{Professorship Big Geospatial Data Management}
\Semester[WiTe 2024/25]{Winter Term 2024/25}
\Vorlesung{Principles of Programming}
\InstitutsWebsite{https://www.bgd.lrg.tum.de/}
\VorlesungsWebsite{https://www.moodle.tum.de/course/view.php?id=103704}

\nummerUebungsblatt{8}
\nummerErsteFrage{1}

\LearningOutcome{ Solving complex computational problems}

\begin{document}
\maketitle

\vspace{10pt}

In this longer tutorial, we want to jointly look at how we can combine our initial knowledge to build a full deep learning system in which a large spatial image and GDAL are used together with Tensorflow to predict a certain spatial feature.

\vspace{10pt}

In this second part, we will learn about Keras data loaders and train a neural network to detect some spatial features.
It is proposed that knowledge on Keras be first built with a tutorial for a simple deep-learning example. 
A beginner-friendly overview can be found \href{https://www.tensorflow.org/tutorials/keras/classification}{here}.
Note that for multidimensional raster data, you may want to use convolutional layers like shown \href{https://keras.io/examples/vision/mnist_convnet/}{here}.
You may even extend the example with RGB images already (you could build an RGB image from three identical grayscale layers).


\begin{task}{Data loading}{}{}
A convenient way to load and prepare your data for deep learning with Keras is a "data loader".
You can read the official documentation and examples \href{https://keras.io/api/data_loading/}{here} and subsequently \href{https://keras.io/api/data_loading/image/}{here}.
Set up data loaders for train, test, and validation data.
\end{task}

\begin{task}{Deep learning model}{}{}
Build and compile a deep learning model with Keras.
You may adapt the convolutional neural network from \href{https://keras.io/examples/vision/mnist_convnet/}{this} example.
\end{task}

\begin{task}{Training and evaluation}{}{}
Finally, train (\code{model.fit()}) and evaluate (\code{model.evaluate()}) your model.
Plot some examples for correct and wrong classified images.
\end{task}

\clearpage

\begin{task}{Improvements}{}{}
Think about ways to improve your model.
You may, for example, try to use data augmentation (which is quite easy with Keras data loaders) or enlarge your training area (and thus the number of training samples).
Further, you may increase or decrease the learning rate and the number of training epochs.

\end{task}

\end{document}